\section{重建不是回到原处}

战后的省份首先面对的不是一个抽象的“恢复”，而是河道、仓储、道路和城镇能否再次把人、粮和消息接在一起。地方官要呈报灾情和军需，商人、船户与地方组织要判断哪一条路线仍可通行，离散者则在征发、借贷、避难和返乡之间寻找余地。朝廷能够发出命令，并不表示银粮已经抵达，或一座城周围的市场已经重新运转。省级军政网络正是在这些断裂处获得作用，也把筹饷、运输和治安的成本留在地方社会。

这一章所能追踪的，是几个危机节点如何改变了重建的条件，而不是为所有地区写出同一条复原曲线。实录、军机档和督抚呈报较能保存朝廷获知何事、怎样下令以及哪些资源被要求调动；它们对远处道路是否畅通、征购怎样落到家庭、谁能离开或留下，常常说得很少。地方志、账簿、契约、日记和报刊可以补出一些地点的断面，却又受保存、书写者和城市视野限制。于是，下文把文书中的行政语言同河流、城镇和边路的实际约束分开，也不把任何一方报告中的规模、损失或控制范围当作不待检验的结论。

\subsection{长江失去一个枢纽后，省份怎样筹措}

\eventanchor{EQ-1853-NANJING-TAKING}{太平军攻占南京并改名天京（1853）}，使长江上一座大城成为长期政权中心。城名的改变并不只是一项象征：江面通行、周边府县的粮食往来、商路的风险和守城所需的物资，都要随之重新安排。对清廷而言，原先可以借助的江南财政与文化枢纽已成为必须争夺和绕开的空间；对沿江居民而言，谁在何处征发、谁能保护船只和货物，往往比远处的政权名称更直接。

朝廷的战事记录和地方呈报能确认南京易手及其后的军事压力，却不能替每一段江路说明货物是否照常流动，或每一处乡村如何筹措军粮。把“长江税源”写成一个可以随城池一并转移的整体，反而会抹去口岸、渡口、乡镇和省界之间的差异。战时筹饷之所以越来越依赖地方中介，不是因为某一道命令突然废除了中央财政，而是因为征解、转运和验收都要穿过不稳定的水陆路径。

\subsection{一条新河道与一个移动战场}

\eventanchor{EQ-1855-YELLOW-RIVER-COURSE-SHIFT}{铜瓦厢决口、黄河改道（1855）}把华北的堤防、耕地、聚落和运输条件推入新的不确定性。决口及河道北流是可以定位的节点；至于不同河段何时受淹、迁徙持续多久、粮价怎样变化，则不能从一份奏报推及整条黄河。河工文书记录的是修防、请款和呈报的行政过程，不等于堤防即时奏效；地方材料即使留下灾情，也往往只照见保存下来的县份和书写者。

同年，\eventanchor{EQ-1855-NIAN-REBELLION-EXPANSION}{捻军活动扩大（1855）}，又使淮北、安徽和华北的村落、堤防与商道进入另一种战争节奏。灾荒、流民、盐业与地方武装网络为行动提供了条件，但不能把水患直接写成一支武装的单一原因。骑兵、团练、运粮者和避难者各自沿着可走的路寻找资源或安全；他们的路线有时相交，有时彼此阻断。官府以“匪”概括的对象并不因此成为一个目的相同的群体，地方材料所能保留的，也只是这些选择的零散痕迹。

两件发生在同一年的事没有一条自动相连的因果线，却共同压缩了财政和运输的余地。河道变化使既有水路与堤工需要重新估量，移动战场又使米粮、牲畜和劳力更难按原有路径调度。省份此后不是被动承受中央压力的容器：地方官、绅商、团练与居民在可用道路和有限信用之间作出的安排，构成了清廷此后能够继续行动的条件，同时也使负担的去向更不均衡。

\subsection{天京的裂口与地方网络的代价}

\eventanchor{EQ-1856-TIANJING-INCIDENT}{天京事变（1856）}在太平天国核心内部造成权力冲突和杀戮，削弱了其协调军政资源的能力。现存的清方记录、奏折与后出的叙述能够提示这一裂口的存在及其政治后果，却不能为城内每一项行动、每一位死者或每一种动机给出无争议的数目和解释。将它写作一场立即决定所有战局的宫廷戏剧，同样会遮住沿江军队、乡村和对手在消息迟缓、补给不足中作出的不同判断。

裂口之所以与省份重建有关，在于军队和行政并不会因首都内部的冲突而停止索取粮饷。谁能筹款、募勇、保住一段航道或维持一座城周边的交换，便在战局中取得新的位置；但这种能力也以地方借贷、捐输、征购和长期驻军为代价。后来的省级军政网络因而既是重建的工具，也是把战争遗留的财政义务和武装关系嵌入地方的渠道。它解决的是协调资源的迫切问题，不能据此推断所有地区已经恢复秩序。

\subsection{京师危机所暴露的运输上限}

\eventanchor{EQ-1860-BEIJING-CAMPAIGN}{英法联军进逼北京（1860）}时，天津一线的谈判、京师防务和热河方向的行幸在短期内相互挤压。圆明园遭焚掠与咸丰帝出走是清楚的节点；英法、清廷和地方观察者留下的记录却服务于不同的军事、外交和政治目的。它们能够说明危机怎样被陈述和处理，不能仅凭一方的叙事便量定各处兵力、开支、掠夺范围或地方反应。

北京的危机也不应被写成压倒内地一切经验的单一中心事件。它与黄河改道、长江战场之间没有共同的即时因果，却都迫使清廷在有限的船路、驿递、饷项和官员注意力中作选择。海路和天津的交涉增加了信息与外交的压力，内战地区则继续要求粮秣和转运。能够抵达北京的呈报，未必等于同一消息已经抵达沿线州县；能够谈定的文本，也不等于港口、道路或省级财政已按其条款改变。重建于是始终受制于运输与信息的速度。

\subsection{西北不是内地战局的余波}

\eventanchor{EQ-1865-DUNGAN-REVOLT}{陕西、甘肃等地战事扩展（1865）}后，西北的城镇、商路、宗教社群和地方军队被卷入多年冲突。地方冲突、军政失序和官府的族群分类相互激化，是理解战争条件的起点，却不足以把不同地方的行动者归入一个单一阵营。官书中的“回”等称谓首先属于行政和征剿的语言；它不能替代人们的自称、宗教联系、市场位置或在战争中的具体选择。

这里的边地不是等候中央重新填补的空白。陕西、甘肃与通往更西方的道路，连接着城镇、绿洲、牧地和跨境贸易的多种节奏；一段道路的失守或重开，会分别影响行军、粮运、逃亡和商业，却未必在所有地点产生同样结果。清廷的军机文书和战事汇编能够追踪调兵、命令和清方所声称的进展；要判断社区边界如何变化、市场如何中断、迁徙和暴力留下何种后果，仍须并读地方材料以及满、蒙、维吾尔、俄文和其他能够取得的区域材料。许多这样的材料保存不全、译名未必相合，正是我们不能把西北战争压缩成一张中央战报的原因。

从省份财政看，西北战争再次显出距离的成本。饷项不是一个从京师发出便自然抵达前线的数目，而是一连串筹措、折算、押运、采购和交接；任何一个环节受阻，前线和沿线社会承受的压力都会改变。现有材料足以说明长期战争要求跨省调度，不能据此替每一处征发、每一批粮饷或每一条边路断言执行的程度。所谓重建，在这里首先是重新取得有限的通行和供给能力，而非对地方生活的即时掌握。

\subsection{天津：地方冲突如何进入条约网络}

\eventanchor{EQ-1870-TIANJIN-MASSACRE}{天津教案（1870）}让育婴堂、传教机构、地方传闻、官府处置和法国外交照会在同一座城市相互碰撞。法国领事、修女和中国民众的死伤使冲突迅速进入外交交涉，但不能因此把天津居民、教民、传教士或官员各自写成没有分歧的阵营。谁相信何种传闻、谁试图调停、谁在暴力中失去行动空间，须回到文本产生的地点和目的中辨认。

\begin{event}{天津教案：地方传闻何以变成外交危机}{同治九年（一八七〇年）}{天津民众、天主教机构、清廷与法国}{冲突及外交后果可核；死伤、谣言传播和个别行动者动机在中外材料中差异很大}
总理衙门档、地方呈报与法国外交材料可以互校危机如何被报告和追责。照会所能确证的是国家怎样提出责任与要求；地方档案和教会通信则分别保留城市社区与宗教网络的片段。它们对人数、责任与意图并不一致，也都不能自行代替那些没有留下文字的人。

天津使战后重建面对另一种限制：省级官员处理地方治安时，已经必须估量条约关系、领事保护和跨国信息的后果。可是，外交文本不能反过来证明城市秩序已经被外来制度接管。保留这种不一致，才能看见教案既是地方社会的撕裂，也是交通、报信和条约网络如何改变行政选择的事件。
\end{event}

\subsection{制度得失：能调动，不等于能复原}

从南京到天津、从黄河到西北，清廷重新取得行动能力的方式往往是借助省级筹饷、地方军政和跨区域运输。这样的安排能在危机中把命令、物资和人手连接起来，却把借贷、征购、驻军和道路风险分配给不同地区与人群。它并未把帝国恢复为战前的样子，而是形成了一套必须依靠地方节点、港口信息和边路通行才能运作的秩序。

这也是证据最容易诱发过度判断之处。制度章程、战报或外交照会能让我们看见计划、要求和国家间的措辞；地方执行、实际产出和普通人的经验需要另一类材料，且常常无法完整取得。读者应把省份的“重建”理解为持续的、地域不均的协商过程：有些路被重新接上，有些社区仍在战争和迁徙的余波中，而边地的行动者从未只是中央叙事的背景。
